\section{Other GPU kernels}
\label{sec:othergpu}

The ECC2K-130 client is one of several GPU arithmetic libraries in
the two repositories. The others target the 256-bit prime field of
secp256k1 (a Koblitz curve in the SEC sense, with $j=0$ and the
order-3 endomorphism), the NIST primes and the polynomial-basis field of
ECC2K-95. They follow the same discipline as the
ECC2K-130 client: a Python or Rust oracle, a host emulation of the
identical device source, bit-identical steppers, planted logarithms,
and static instruction accounting from the assembler. What most of
them still lack is a measurement on a device. We start by saying
which is which, because a static instruction count and a measured
rate are different kinds of evidence, and the rest of the paper keeps
them apart.

\subsection{Status}
\label{sec:othergpu-status}

\begin{table}[t]
\centering
\small
\caption{GPU arithmetic in the two repositories by evidence class.
  ``Emulated'' means the device source was executed on the host and
  checked against an oracle; ``static'' means instruction, register
  and stack counts from \texttt{ptxas} with no execution.}
\label{tab:gpu-status}
\begin{tabular}{@{}p{0.30\textwidth}p{0.17\textwidth}p{0.22\textwidth}p{0.22\textwidth}@{}}
\toprule
kernel family & field & on a device & otherwise \\
\midrule
ECC2K-130 client (Sections~\ref{sec:gpu}, \ref{sec:campaign}) & $\FF_{2^{131}}$, ONB & measured, audited & emulated, static \\
secp256k1 / Curve25519 rho, baby-step giant-step & $\FF_p$, two 256-bit primes & never & emulated, static \\
secp256k1 kangaroo & $\FF_p$, 256-bit & never & emulated, static \\
P-256 / P-384 field and rho walk (Section~\ref{sec:nist}) & $\FF_p$, 256- and 384-bit & measured (RTX PRO 6000, 2026-10-12) & host-verified; protocol pre-registered \\
ECC2K-95 kernel & $\FF_{2^{97}}$, polynomial basis & never & emulated, static \\
64-bit prime-field rho (\texttt{cryptanalysis}) & $\FF_p$, $p<2^{64}$ & never & emulated, static \\
\bottomrule
\end{tabular}
\end{table}

Table~\ref{tab:gpu-status} is the inventory. Every ``never'' repeats a
sentence in that family's own README, and the reason is the same each
time: the code was written and verified in containers with no GPU or
no complete CUDA toolkit, and the device runs were planned for rented
hardware that the campaign's account hold
(Section~\ref{sec:campaign-state}) or missing credentials then put out
of reach. The NIST kernels were in the same position until they
finally ran, which is the subject of Section~\ref{sec:nist}~\evcite{\evref{ecc}; \evref{kangaroo}; \evref{ecc2k95}; \evref{nistattempt}; \evref{prho64}, ``How this is tested''}.
The numbers below are therefore of two kinds, and we label each one: a
static count, or an emulated operation count or walk constant.

\subsection{secp256k1}
\label{sec:othergpu-fp256}

The field layer holds an element in eight 32-bit limbs, because the
integer datapath is 32 bits wide and the carry-chain instructions
\texttt{mad.lo.cc}, \texttt{madc.hi.cc} and \texttt{addc} exist only at
that granularity. For secp256k1 the 512-bit product is folded with
$2^{256}\equiv2^{32}+977 \pmod p$, about $40\%$ fewer multiply-adds
than Montgomery; for any other odd modulus below $2^{256}$ a
Montgomery CIOS path is used, and the scalar ring modulo $n$ reuses
it on the host. Points are Jacobian with the standard $a=0$ doubling
($2M+5S$) and mixed addition ($7M+4S$) formulas. The library is a
throughput library, not a signing library, and is not
constant-time~\evcite{\evref{ecc}; \evref{ecc}}.
A second special-reduction path, added in October 2026, serves
Curve25519 (Ed25519's prime-order group in short-Weierstrass form,
$p=2^{255}-19$, folded with $2^{256}\equiv38$) through the same
walk, baby-step giant-step and test suites; it is host-verified
against the RFC~7748 constants and has not run on a
device~\evcite{\evref{ecc}}.

\begin{table}[t]
\centering
\small
\caption{Static PTX instruction counts for one 256-bit field
  multiplication (clang, \texttt{sm\_90}, \texttt{-O3}). The hard floor
  for the product is 128 instructions: 64 limb products at two
  instructions each. Class: engineering; nothing here has executed on
  a device.}
\label{tab:fp256}
\begin{tabular}{@{}lrrr@{}}
\toprule
 & portable C++ & inline PTX, before & inline PTX, after \\
\midrule
512-bit product & 236 & 171 & 137 \\
secp256k1 reduction & 146 & 130 & 77 \\
full multiply & 472 & 328 & \textbf{210} \\
\midrule
Montgomery CIOS multiply & 475 & --- & 352 \\
\bottomrule
\end{tabular}
\end{table}

Table~\ref{tab:fp256} is the instruction ledger. The portable path
emits no carry instructions at all: of the 444 instructions in a
portable squaring, 178 are adds, 102 shifts, 100 ANDs and only 46
multiplies, so the multiply is carry-bound, not multiply-bound, and
the inline-PTX rewrite's $55\%$ cut comes from the carry chains. The
seven assembly blocks are checked by an interpreter against the
portable path on $20{,}000$ random and boundary inputs each, and that
check is mutation-tested; it does not cover register allocation or
device semantics, which is why \texttt{FP\_PTX} is off by
default~\evcite{\evref{eccopt}, §2; \evref{ecc}}.

In the unit of Section~\ref{sec:arith}, a secp256k1 inversion is 270
multiplications (255 squarings and 15 multiplications along the
libsecp256k1 addition chain), a rho step with batch $W$ is $6+270/W$
multiplications ($3$ for the affine addition, $3$ for Montgomery's
trick), $7+270/W$ with the $j=0$ fold: $39.8$ at $W=8$, $14.4$ at
$W=32$, $8.1$ at $W=128$, against 276 unbatched~\evcite{\evref{eccopt}, §1}.
Three things were tried and declined on the static count: a
dedicated squaring routine ($-2.7\%$ portable, $+23.5\%$ on the PTX
path; reverted), lazy reduction to $[0,2^{256})$ (worth about $10\%$
of a multiply, but it breaks the canonical form that equality, the
distinguished-point test, the cycle hash and the fold comparison rely
on), and Karatsuba (judged to lose on a carry-bound multiply). Tensor
cores were considered and not reached for~\evcite{\evref{eccopt}, §2 and §4.4}.
Marking the inversion and ladder routines \texttt{\_\_noinline\_\_}
and replacing two 16-entry tables by a small double-scalar ladder
took the walk kernel's stack from $4{,}664$ to $1{,}616$~B (the
low-memory variant from $4{,}008$ to $792$~B) and the compile from
over ten minutes to 46~s, the same code-size lesson as
Section~\ref{sec:arith-bitsliced}~\evcite{\evref{eccopt}, §3}.

\subsection{Automorphism folding and BSGS}
\label{sec:othergpu-rho}

The rho walk is $r$-adding on the low bits of $x$, with distinguished
points by a mask on the next byte and three folds: identity, the
negation map, and on $j=0$ curves the full $\langle\beta,-1\rangle$
orbit of order six, chosen by minimising $x$ over
$\{x,\beta x,\beta^2 x\}$ at the cost of one multiplication
($\beta^2x=-(x+\beta x)$) and two compares, about $10\%$ of a batched
step for $\sqrt3$ fewer steps. Only an automorphism code is kept per
step; a report is $\{x,\text{walk},\text{restart},\text{steps}\}$, 48
bytes, and the host replays the two colliding walks from their seeds,
multiplying $(a,b)$ by $(-1)^{\text{neg}}\lambda^k$ at each step. Fruitless
cycles of length 2 to 6 are detected with a five-deep ring and escaped
by doubling the member with the smallest hash; their measured rates
at $R=8$ agree with the prediction to about
$10\%$~\evcite{\evref{ecc}, ``The walk''}.

On a 36-bit $j=0$ toy curve in host emulation (8 solves per fold, 16
walks, $R=256$, period $2^7$), the mean steps were $305{,}740$,
$199{,}042$ and $113{,}522$ for folds 1, 2 and 6 against predictions of
$249{,}454$, $176{,}391$ and $101{,}839$ ($1.23\times$, $1.13\times$,
$1.11\times$; $\pm0.18\times$ on a mean of eight), all $8/8$ correct,
so fold 6 buys $1.75\times$ over the negation map (predicted $1.73$).
In the unit $S=\text{ops}/\sqrt n$ that is $1.54$, $1.00$ and $0.57$.
The three steppers (batched, low-memory, one-inversion-per-walk) are
bit-identical over hundreds of iterations with the same
distinguished-point multiset, and every point replays on the host to
$aP+bQ$, which is what proves the $\lambda$ bookkeeping on secp256k1
itself~\evcite{\evref{ecc}, ``How this is tested''}.

Baby-step giant-step on the same arithmetic keeps an $x$-keyed table
of 16-byte entries in lock-free open addressing (one 64-bit
\texttt{atomicCAS} per probe, about $2.5$ slots per probe at half
load, a false-candidate rate of load$/2^{32}$ and none observed), and
walks $W$ giant chains per thread sharing one inversion, 94 registers
and no spills at $W=8$. With the negation map the layout is $m=\sqrt{w}/2$
babies and $2m-1$ giants, the Galbraith--Wang--Zhang row at $1.00\sqrt w$
expected against the textbook $1.41\sqrt w$. A $2^{66}$-wide interval
is $2^{32}$ entries and 64~GB, the top of one 80~GB device; beyond
that the method is out of memory, not out of
time~\evcite{\evref{ecc}, ``Baby-step giant-step''}.
Table~\ref{tab:bsgs} is the measured cost on a 40-bit toy curve, in
emulation, every phase charged.

\begin{table}[t]
\centering
\small
\caption{Rho and baby-step giant-step on a 40-bit toy curve
  ($n=649{,}523{,}094{,}257$), host emulation of the device source, 128
  chains, mean over 8 random targets. $S=\text{group additions}/\sqrt n$
  with table construction charged. The rho rows predate the order-six
  fold and were single solves ($\pm0.5\times$). Twelve targets on a
  $2^{26}$ interval cost $S=0.560$ per target against $2.131$ for one
  cold solve.}
\label{tab:bsgs}
\begin{tabular}{@{}lrrrl@{}}
\toprule
row & table$/\sqrt n$ & mean $S$ & predicted & correct \\
\midrule
rho, negation map (reference) & 0 & 0.85 & --- & 1/1 \\
rho, no negation map & 0 & 1.42 & --- & 1/1 \\
BSGS textbook, reference stepper & 0.707 & 1.44 & 1.41 & 21/21 \\
BSGS textbook, batched $W=16$ & 0.707 & 1.44 (bit-identical) & 1.41 & 21/21 \\
BSGS negation layout, $W=16$ & 0.500 & 1.07 & 1.00 & 21/21 \\
\quad giant phase only & 0.500 & 0.56 & 0.50 & 8/8 \\
\bottomrule
\end{tabular}
\end{table}

The classification is the same as every row in Section~\ref{sec:gpu}:
the batched stepper is engineering, the layouts are the published
ones, and nothing moves the ratio to the generic floor. A Rust port of
the same baby-step giant-step on the CPU runs at $37$--$49$~M
steps/s on four threads at $S=1.03$ cold and $0.55$ amortised, $6$--$11\times$
the big-integer implementation it replaced~\evcite{\evref{ecc}, ``Cost, measured''}.
The emulation itself runs at about $380$~k steps/s per core, a
figure that says nothing about a device and is quoted only so that
no one mistakes it for one.

\subsection{Kangaroo}
\label{sec:othergpu-kangaroo}

The interval solver reuses the verified field and point code for
$Q=kG$ with $k\in[a,a+w)$ at about $2\sqrt w$ steps: two herds that
alternate by thread parity so every warp is half tame and half wild,
a jump table of $2^6$ pseudo-random scalars with mean jump $\sqrt w/2$
and 25-word entries (odd, so the shared-memory reads are
bank-conflict-free), and the target centred at $Q'=Q-(a+w/2)G$. The
invariant that makes it testable is that after every step a tame
kangaroo at distance $d$ is at $dG$ and a wild one at $Q'+dG$, which
pins the jump table, the distance accumulator and the point
arithmetic together~\evcite{\evref{kangaroo}, ``How the walk works''}.

Its measured constant is the useful number. In emulation on a $2^{31}$
interval with 120 trials per configuration, steps are $C\sqrt w$ with
$C=2.79\pm0.13$ for the default (median $2.65$); the idealised
two-herd bound is $2.0$, Pollard's classic lambda about $3.3$, and
published parallel implementations about $2.1$, so this one is about a
third above the best published. A jump-parameter sweep at 20 trials
per cell had shown $C$ from $2.50$ to $3.95$, an apparent $36\%$ win;
rerun at 120 trials the best cell and the default were within half a
standard error. The defaults sit on a flat optimum and the sweep was
noise, which is recorded so that it is not rediscovered. Restarting
at a distinguished point or not was likewise within
noise~\evcite{\evref{kangaroo}, ``Measured cost''}. A multi-device
split (one host thread per device, disjoint even index bases so herd
parity stays global, one host table) is tested by reproducing the
combined herd exactly from two contexts and recovering a 28-bit key
across them; multi-machine is not done~\evcite{\evref{kangaroo}, ``Multi-GPU''}.
The tool ships no real puzzle keys (no trustworthy offline source),
does no address scanning, and has not solved any puzzle; the
synthetic 20-, 28- and 32-bit registry entries are regression fixtures.

\subsection{NIST curves}
\label{sec:othergpu-nist}

The NIST library has graduated from this inventory: it now has device
measurements, and Section~\ref{sec:nist} reports them as the ECC2K-130
ladder carried over to prime fields. One historical note belongs
here. The library's first device attempt was an infrastructure
non-result, because the cloud credentials were not exposed to the CI
runner and the GPU steps were skipped. The attempt record says it
``must not be cited as performance evidence'', and we do not cite
it~\evcite{\evref{nistattempt}}.

\subsection{Other binary fields}
\label{sec:othergpu-binary}

\paragraph{The ECC2K-95 kernel.} Section~\ref{sec:ecc2k95} used this
kernel as the end-to-end control; its arithmetic is the packed,
polynomial-basis counterpart of Section~\ref{sec:arith-packed}. An
element is four 32-bit words; a 32-bit carry-less product is sixteen
widening integer multiplies with the interleaved-mask trick (operands
split by bit index modulo 4 so no partial sum exceeds 8 and nothing
carries); the 1-bit top word at $m=97$ is split off so that a
three-way Karatsuba over three words (six leaves) plus seven
single-bit masks replaces nine leaves; reduction by $t^{97}+t^6+1$ is
two folds, the second an extract and two XORs; squaring is a
bit-spread kept in 64-bit form (78 instructions against 104 for two
32-bit spreads); inversion is Itoh--Tsujii at about 7 multiplications
and 96 squarings. The static count is 304 instructions per
multiplication, down from 405, with the six software leaves 96 of the
304; the native \texttt{clmad} that gave the ECC2K-130 client $+22\%$
would do the same product in six instructions and is not yet ported
because the pip-wheel toolchain stops at \texttt{ptxas}
12.9~\evcite{\evref{ecc2k95}, ``What the field arithmetic actually costs''; \evref{ecc2k95}}.

Per step the cost model is one inversion split $W$ ways plus two
multiplications, one squaring and the three multiplications of
Montgomery's trick, with the two Frobenius chains of the walk paid in
squarings: $3{,}813$ instructions at $W=8$, $44\%$ multiplies, $27\%$
the Frobenius chains, $25\%$ the inversion's squarings. Because
$\tau^k$ is $\FF_2$-linear it can be applied as 25 four-bit-window
table reads at any $k$, 296 instructions against 78 per squaring;
tables for $\tau^{6,12,24,48}$ cut the inversion's $3{,}744$
squaring instructions at $k=48$ to 296 and the step to $3{,}083$
($-19\%$), and $W=4$ with tables beats $W=8$ without at 656 against
$1{,}048$ bytes of local memory. The table traffic, 400 words per
inversion, is unmeasured, and the tables are off by
default~\evcite{\evref{ecc2k95}, ``$\tau^k$ is not $k$ squarings''}.
The low-memory stepper at $W=8$ compiles to 128 registers and 464~B
of stack with four blocks per SM, 512 threads per SM on
\texttt{sm\_90}, \texttt{sm\_100} and \texttt{sm\_120}; the
prime-field walk of Section~\ref{sec:othergpu-fp256} needs 226
registers and 792~B at 256 threads per SM, so the binary walk gets
$1.5$--$2\times$ the occupancy before either has run~\evcite{\evref{ecc2k95}, ``Occupancy''}.
A 40-bit end-to-end test of eight logarithms on one CPU core fell
from $2.1$--$2.7$~s to $1.02$--$1.05$~s across the rewrite at the same
$191{,}840$ steps, which the README itself files as engineering with
$S$ flat.

\subsection{64-bit prime-field rho}
\label{sec:othergpu-prho64}

The \texttt{cryptanalysis} library's 64-bit prime-field CUDA rho (one limb, Montgomery
with $R=2^{64}$, eight walks per thread sharing a Fermat inversion so
that a step is 18 multiplications instead of 99) is in the same
state as most kernels above: $77{,}000$ device-arithmetic assertions
against the host library, whole-solver emulation on four groups, a
cubin for \texttt{sm\_70} through \texttt{sm\_90} with 64--72 registers
and no spills, and no execution on a device; its emulator rows sit at
$S\approx1.3$--$2.6$ alongside the CPU rho and show only that $S$ does
not grow with the walk count~\evcite{\evref{prho64}; \evref{prho64}, ``GPU rho kernel''}.
Its efficiency scorer, built for the ECC2K-130 walk, is the one place
the two threads meet: charging the packed walk's $4{,}328$
lane-instructions per iteration to the measured dual-issue rate of
the RTX~PRO~6000 ($78.3$ per SM-clock, 188 SMs, $2{,}307$~MHz) gives
an attainable $7.85$~billion iterations/s per card against an audited $6.905$ for that
build, $88\%$ of the issue roof, and the same method on the wide
multiply reproduces the ``31~billion iterations/s if the walk did nothing else''
of the earlier ceiling note~\evcite{\evref{gpueff}; \evref{gpueff}}.

\subsection{Summary}

Two lessons hold across the fields. The first is that, on this
hardware, the carry-less unit is the whole difference between the
binary and prime-field kernels. A 131-bit binary product is six
\texttt{CLMAD}, while a 256-bit prime-field product is at best 210
carry-chained integer instructions. The binary walk is therefore
limited by a narrow unit at $1.6$--$2.0$ lanes per SM-clock, the
prime-field walk by carries on the wide integer pipe, and the two want
different occupancies, batch sizes and arithmetic representations.
Section~\ref{sec:nist} confirms this on a device: the same levers
helped both walks, but by very different amounts. The second lesson is
that each of these libraries reached a verified, bit-identical,
statically accounted state without ever touching a device, and that
state is worth exactly what Section~\ref{sec:floor} says a pre-tuning
floor is worth. It bounds what a device could do; it does not tell you
what one did. The NIST walk is a reminder of how wide that gap can be:
its first device run measured an empty kernel. Only the ECC2K-130
client and, more recently, the NIST walk have crossed to the measured
side, which is why they get sections of their own and the rest share
this one.
